% ================================================================
% TAILORING
% ================================================================

\section{Tailoring}
\label{sec:tailoring}

Lo standard ISO/IEC~12207:1995 prevede la possibilità di adattare
(\textit{tailoring}) i processi del ciclo di vita software al
contesto specifico del progetto.

Il gruppo \textit{Synergo Unit} adotta un approccio iterativo orientato al
miglioramento continuo introducendo progressivamente processi e
pratiche in funzione dell'evoluzione del progetto e delle baseline\textsubscript{G} previste
dal Regolamento del Progetto Didattico. Il ciclo di controllo periodico
(consuntivo di periodo, retrospettiva, misure correttive, aggiornamento
della pianificazione) tramite cui il gruppo rivede sia le stime future sia
le decisioni pregresse è descritto in dettaglio in Sezione~\ref{sec:sviluppo}.

Il tailoring adottato tiene conto dei seguenti fattori contestuali:

\begin{itemize}
    \item \textbf{Baseline attuale}: PB\textsubscript{G}; il gruppo ha
    consolidato l'architettura e ha avviato lo sviluppo esteso del prodotto
    (Minimum Viable Product) e della relativa documentazione richiesta per la Product Baseline;

    \item \textbf{Dimensione del gruppo}: sei membri con rotazione periodica e
    documentata dei ruoli;

    \item \textbf{Natura didattica del progetto}: il progetto viene svolto in ambito
    accademico, con vincoli temporali, organizzativi ed economici definiti dal
    regolamento del corso;

    \item \textbf{Introduzione progressiva dei processi}: il gruppo introduce
    esclusivamente processi effettivamente utilizzati, verificabili e coerenti
    con lo stato corrente del progetto;

    \item \textbf{Principio just-in-time}: ogni nuova attività o pratica viene normata
    prima della sua applicazione operativa.
\end{itemize}

Il presente documento è soggetto ad aggiornamento progressivo: ogni evoluzione
significativa del way of working, degli strumenti o delle attività
di progetto comporta, se necessario, l'aggiornamento delle presenti Norme di
Progetto.

% ------------------------------------------------

\subsection{Processi Attivati nella Fase CC}

La seguente tabella riporta i processi attivati durante la fase di
candidatura\textsubscript{G} (CC), mantenuti operativi anche nelle fasi successive
del progetto.

\bigskip

\begin{xltabular}{\textwidth}{
    >{\raggedright\arraybackslash}p{3cm}
    >{\raggedright\arraybackslash}X
    >{\raggedright\arraybackslash}p{3cm}
}
    \toprule
    \textbf{Categoria} & \textbf{Processo} & \textbf{Sezione} \\
    \midrule
    \endfirsthead

    \toprule
    \textbf{Categoria} & \textbf{Processo} & \textbf{Sezione} \\
    \midrule
    \endhead

    \midrule
    \multicolumn{3}{r}{\textit{Continua nella pagina successiva}} \\
    \endfoot

    \bottomrule
    \endlastfoot

    Primario
    & Processo di Fornitura
    & \ref{sec:fornitura} \\

    Di supporto
    & Processo di Documentazione
    & \ref{sec:documentazione} \\

    Di supporto
    & Processo di Gestione della Configurazione
    & \ref{sec:configurazione} \\

    Di supporto
    & Processo di Verifica
    & \ref{sec:verifica} \\

    Organizzativo
    & Processo di Gestione
    & \ref{sec:gestione} \\

    Organizzativo
    & Processo di Comunicazione
    & \ref{sec:comunicazione} \\

    Organizzativo
    & Processo di Infrastruttura
    & \ref{sec:infrastruttura} \\

\end{xltabular}

% ------------------------------------------------

\subsection{Processi Attivati nella Fase RTB}

Con l'avanzamento alla baseline RTB vengono introdotti ulteriori processi e
attività necessari al consolidamento dei requisiti, alla motivazione delle scelte
tecnologiche e all'implementazione dimostrativa del Proof of Concept\textsubscript{G}.

Il Proof of Concept è finalizzato a dimostrare la fattibilità tecnica delle
funzionalità principali individuate per la Requirements \& Technology
Baseline; non costituisce ancora il prodotto completo previsto
per la Product Baseline\textsubscript{G}.

I processi precedentemente attivati rimangono operativi.

\bigskip

\begin{xltabular}{\textwidth}{
    >{\raggedright\arraybackslash}p{3cm}
    >{\raggedright\arraybackslash}X
    >{\raggedright\arraybackslash}p{3cm}
}
    \toprule
    \textbf{Categoria} & \textbf{Processo} & \textbf{Sezione} \\
    \midrule
    \endfirsthead

    \toprule
    \textbf{Categoria} & \textbf{Processo} & \textbf{Sezione} \\
    \midrule
    \endhead

    \midrule
    \multicolumn{3}{r}{\textit{Continua nella pagina successiva}} \\
    \endfoot

    \bottomrule
    \endlastfoot

    Primario
    & Processo di Sviluppo (Analisi dei Requisiti, Analisi dello Stato dell'Arte, Progettazione Architetturale Preliminare, Implementazione del Proof of Concept)
    & \ref{sec:sviluppo} \\

    Di supporto
    & Processo di Validazione\textsubscript{G}
    & \ref{sec:validazione} \\

    Di supporto
    & Processo di Accertamento della Qualità
    & \ref{sec:qualita} \\

    Di supporto
    & Processo di Gestione dei Cambiamenti
    & \ref{sec:cambiamenti} \\

    Di supporto
    & Processo di Gestione delle Non Conformità\textsubscript{G}
    & \ref{sec:non-conformita} \\

    Organizzativo
    & Gestione dei Rischi
    & \ref{sec:gestione} \\

    Organizzativo
    & Processo di Miglioramento
    & \ref{sec:miglioramento} \\

    Organizzativo
    & Processo di Formazione
    & \ref{sec:formazione} \\

\end{xltabular}

% ------------------------------------------------

\subsection{Processi Attivi nella Fase PB}

Con il superamento della RTB e l'ingresso nella Product Baseline, il gruppo ha consolidato l'architettura e avviato lo sviluppo esteso del prodotto (Minimum Viable Product).

I processi precedentemente attivati rimangono operativi. La seguente tabella riporta i processi formalmente normati e introdotti in questa fase per supportare lo sviluppo del codice di produzione; il dettaglio è riportato nei corrispondenti capitoli dei Processi Primari, Organizzativi e di Supporto.

\bigskip

\begin{xltabular}{\textwidth}{
    >{\raggedright\arraybackslash}p{3cm}
    >{\raggedright\arraybackslash}X
    >{\raggedright\arraybackslash}p{3cm}
}
    \toprule
    \textbf{Categoria} & \textbf{Processo} & \textbf{Sezione} \\
    \midrule
    \endfirsthead

    \toprule
    \textbf{Categoria} & \textbf{Processo} & \textbf{Sezione} \\
    \midrule
    \endhead

    \midrule
    \multicolumn{3}{r}{\textit{Continua nella pagina successiva}} \\
    \endfoot

    \bottomrule
    \endlastfoot

    Primario
    & Processo di Sviluppo (Specifica Tecnica, Manuale Utente, codifica estesa del prodotto, definizione dell'architettura, norme di codifica e naming conventions, Dependency Injection)
    & \ref{sec:sviluppo} \\

    Di supporto
    & Processo di Gestione della Configurazione (gestione rigorosa delle Pull Request)
    & \ref{sec:configurazione} \\

    Di supporto
    & Processo di Verifica (test automatici Unit e Integration testing, analisi statica del codice, Continuous Integration tramite GitHub Actions)
    & \ref{sec:verifica} \\

\end{xltabular}